\documentclass[11pt]{article}
\usepackage[margin=1in]{geometry}
\usepackage{amsmath, amssymb, amsthm}
\usepackage{hyperref}
\usepackage{booktabs}

% ============================================================
% preprint_draft_v4.tex
% Supersedes preprint_draft_v3.tex.  Changes in v4, all forced by the
% 2026-07-25 certification round:
%
%  1. HEADLINE 0.3802973 -> 0.3803953504.  The old number rested on an
%     anchor convention (solver value minus 1e-5), not on a certificate.
%     Every anchor now carries a Jansson-Chaykin-Keil interval-arithmetic
%     bound, and the certified anchors came out ABOVE the convention, so
%     replacing the convention by a theorem RAISED the bound.
%
%  2. PRIORITY CLAIM REMOVED.  v3 said "the first lower-bound advance on
%     this problem since 2022".  That is false: Kim & Pilanci (ICML 2026)
%     certified 0.37912 in June 2026.  We now lead the published record by
%     1.28e-3, which is the correct and still-strong claim.
%
%  3. SATURATION THEOREM WITHDRAWN.  v3's full-stack ceiling C_explicit =
%     0.380713 is not sound as derived: its own per-row table yields
%     0.381586 + 5.85e-4 = 0.382171, which exceeds the certified upper
%     bound, and the published figure substitutes a min-over-rows anchor
%     for the theorem's per-row right-hand side.  Replaced by a
%     feasible-set containment statement, which is weaker, correct, and
%     still sufficient to conclude the architecture cannot close the gap.
%
%  4. NEW: certified upper bound, and a correction to the record.
% ============================================================

\newtheorem{theorem}{Theorem}
\newtheorem{proposition}[theorem]{Proposition}
\newtheorem{lemma}[theorem]{Lemma}
\theoremstyle{definition}
\newtheorem{remark}[theorem]{Remark}

\title{A certified lower bound for the Erd\H{o}s minimum overlap constant}
\author{Ben Zanghi}
\date{July 2026}

\begin{document}
\maketitle

\begin{abstract}
Let $\mu$ denote the Erd\H{o}s minimum overlap constant. We prove
$\mu \ge 0.3803953$, improving the best published lower bound
$0.37912$ of Kim and Pilanci by $1.28 \times 10^{-3}$ and White's
$0.379005$ by $1.39 \times 10^{-3}$. The bound augments White's
Fourier-analytic convex program with Bochner moment-matrix and
polynomial-moment constraints over a cover of the residual parameter
region.

The methodological point of this paper is not the augmentation but the
certification. Each of the twelve dual certificates underlying the cover
is verified by a Jansson--Chaykin--Keil a-posteriori bound computed in
directed-rounding interval arithmetic, with the dual multipliers read
from the same solve that produced the certificate; no step trusts the
solver's status flag, its iteration log, or its floating-point
arithmetic. This replaces a widely convenient but unjustified
convention---taking a solver-reported objective less a fixed
margin---and, because the certified values exceed that convention at
every centre, it \emph{raises} the resulting bound rather than costing
anything.

We give the bound in two tiers. The stronger, $\mu \ge 0.3803953$, uses
our own cover on the full parameter space. The more robust,
$\mu \ge 0.380000$, is load-bearing on nothing but the twelve
certificates and White's published Table~2, and we recommend it to
readers who want the shortest possible chain of dependence.

Separately we certify $\mu \le 0.380859056614806899090596051449$ in
exact rational arithmetic from a publicly available witness, and we
correct the record on the upper-bound side: the value $0.380856$ in
current circulation is not an upper bound on $\mu$.
\end{abstract}

\section{Introduction}\label{sec:intro}

Erd\H{o}s asked, for a partition of $\{1,\dots,2n\}$ into two sets $A$
and $B$ of size $n$, how small the largest difference-multiplicity
$\max_k |\{(a,b) \in A\times B : a-b = k\}|$ can be made. Writing $M(n)$
for that minimum, the limit $\mu = \lim_{n\to\infty} M(n)/n$ exists, and
admits the continuous formulation
\begin{equation}\label{eq:mu}
  \mu \;=\; \inf_{h}\ \sup_{k \in \mathbb{R}}\
  \int_0^2 h(x)\bigl(1-h(x+k)\bigr)\,dx,
\end{equation}
the infimum over measurable $h:[0,2]\to[0,1]$ with $\int h = 1$, both
$h$ and $1-h$ extended by zero off $[0,2]$.

The two sides of \eqref{eq:mu} are attacked by disjoint machinery. Upper
bounds are exhibitions: any admissible $h$ gives one, and the current
record constructions come from large-scale automated search. Lower
bounds are exclusions: they must rule out every $h$ simultaneously, and
so require a certified relaxation. This paper is about the lower side.

\paragraph{State of the art.} White \cite{White2023} proved
$\mu \ge 0.379005$ by a Fourier-analytic convex program with a
divide-and-conquer cover of a three-parameter space. Kim and
Pilanci \cite{KimPilanci2026} improved this to $\mu \ge 0.37912$ using an
agent-assisted search for valid convex relaxations, certified in interval
arithmetic. On the upper side, a sequence of automated searches has
produced constructions clustering near $0.38086$.

\paragraph{Results.}
\begin{theorem}\label{thm:main}
$\mu \ge 0.3803953$.
\end{theorem}
\begin{theorem}\label{thm:robust}
$\mu \ge 0.380000$, using only the twelve certificates of
Section~\ref{sec:cert} together with White's published Table~2.
\end{theorem}
\begin{theorem}\label{thm:ub}
$\mu \le 0.380859056614806899090596051449$, certified in exact rational
arithmetic.
\end{theorem}

Theorem~\ref{thm:robust} deserves comment. It is numerically weaker than
Theorem~\ref{thm:main} but rests on a far shorter chain: our own
contribution enters only through twelve interval-arithmetic
certificates, and everything outside the residual region is quoted from
a peer-reviewed paper. Where the two disagree in spirit---one maximally
strong, one maximally defensible---we would rather a reader take the
second.

\paragraph{What this paper does not claim.} We do not claim to have
closed, or to be able to close, the gap. Section~\ref{sec:ceiling}
records a containment argument showing that this entire family of
relaxations is bounded above by roughly $0.38065$ at the parameters of
the extremal function, some $2.1\times10^{-4}$ short of the certified
upper bound. Section~\ref{sec:negative} records five further approaches
that we ruled out, with proofs. We regard those exclusions as a
substantive part of the contribution: they say where the remaining
difficulty is not.

\section{The augmented program}\label{sec:setup}

White's program tracks a competitor through cell averages of a function
$M$ on a grid of $N$ cells together with the first $T$ Fourier
coefficients $(c_k,d_k)$ of the associated density $f$, subject to a
family of linear and second-order-cone constraints indexed by a
parameter triple $\theta = (E(M), c_1, d_1)$. We add:
\begin{enumerate}
\item \textbf{Bochner moment-matrix PSD.} For $f$ and for $1-f$, the
  Hermitian Toeplitz matrix of Fourier coefficients is positive
  semidefinite; we impose this at truncation level $40$.
\item \textbf{Polynomial-moment positivity.} $\int x^{2k} f \ge 0$ for
  $k \le 10$, expressed in $(c,d)$ through the Fourier expansion of
  $x^{2k}$, with an \emph{analytic} remainder bounding the truncated
  tail. (An earlier version of this paper used a hard cutoff with no
  remainder; that was an error, corrected here and in
  \cite{repo}. See Remark~\ref{rem:tail}.)
\end{enumerate}

\begin{remark}\label{rem:tail}
Truncating a series that appears inside a valid inequality, without an
analytic bound on the discarded tail, silently strengthens the
constraint and inflates the resulting bound. This failure mode retracted
two intermediate results in the development of this work. Every
truncated sum in the final program carries a proved remainder.
\end{remark}

\paragraph{The residual region.} White's divide-and-conquer reduces the
problem to a single box: combining his Table~2 shows that either
$\mu \ge 0.37925$, or the extremal function's parameters satisfy
\begin{equation}\label{eq:516}
  0 \le E(M^*) \le 0.06, \qquad
  0.35 \le c_1^* \le 0.45, \qquad
  -0.02 \le d_1^* \le 0.02 .
\end{equation}
All twelve of our centres lie in \eqref{eq:516}, and it is the region
that binds in Theorem~\ref{thm:main}.

\section{Certification}\label{sec:cert}

This section contains the paper's main methodological content.

\subsection{The problem with solver-reported bounds}

A conic program solved by an interior-point method returns an
approximate primal--dual pair. The dual objective lower-bounds the
primal optimum \emph{only if the dual point is feasible}, and an
interior-point iterate is feasible only up to its residual. It is
therefore not sound to report a dual objective, or a solver's printed
objective less an ad-hoc margin, as a lower bound. Both are widespread.

\subsection{The Jansson bound, with matched multipliers}

Let the canonical form be $\min c^\top x$ subject to $Ax + s = b$,
$s \in \mathcal{K}$, and let $z$ be the approximate conic dual returned
by the solver. Following Jansson, Chaykin and Keil \cite{JCK2007}, for
every primal-feasible $x$,
\begin{equation}\label{eq:jansson}
  c^\top x \;\ge\; -b^\top z \;+\; \underbrace{\textstyle\sum_j
    \min\bigl(0,\lambda_{\min}^{\mathcal{K}_j^*}(z_j)\bigr)\,
    \bar s_j}_{=:\ \mathrm{pen}_{zs}\ \le\ 0}
  \;-\; \underbrace{\textstyle\sum_i |D_i|\,\bar x_i}_{=:\
    \mathrm{pen}_{Dx}\ \ge\ 0},
  \qquad D = c + A^\top z,
\end{equation}
where $\bar x_i$ are a-priori bounds on $|x_i|$ from the model box
constraints and $\bar s_j$ bound the slack sizes. We write $p_{\mathrm{lo}}$
for the right-hand side of \eqref{eq:jansson}, evaluated with every
arithmetic step performed in directed-rounding interval arithmetic, and
with verified enclosures for the per-cone minimum eigenvalues.

The step that makes this usable inside a parameter cover is the
following. Moving $\theta$ within the residual region changes only the
right-hand side $b$ of the four parameter-dependent constraints; $A$ and
$c$ are untouched. Hence
\begin{itemize}
\item $-b(\theta)^\top z = -b(\theta_c)^\top z + \mathrm{shift}_c(\theta)$
  exactly, with $\mathrm{shift}_c$ White's dual-objective transport;
\item $\mathrm{pen}_{Dx}$ is $\theta$-independent, since $D$ does not
  involve $b$ and the $\bar x_i$ are fixed model bounds;
\item $\mathrm{pen}_{zs}$ is $\theta$-independent \emph{when it
  vanishes}, i.e.\ when $z$ lies in the dual cone.
\end{itemize}

\begin{proposition}\label{prop:transport}
If $\mathrm{pen}_{zs} = 0$ at a centre $\theta_c$, then
$\Phi_c(\theta) := p_{\mathrm{lo}} + \mathrm{shift}_c(\theta)$
is a valid lower bound on $\mu$ conditional on the true parameters being
$\theta$, for every $\theta$ in the region.
\end{proposition}

Our driver asserts $\mathrm{pen}_{zs} = 0$ at every centre and refuses to
emit one otherwise. It held at all twelve. Crucially, the multipliers
defining $\mathrm{shift}_c$ are read from the \emph{same} solve that
produced $p_{\mathrm{lo}}$; pairing a certificate from one solve with
multipliers from another is unsound, since $p_{\mathrm{lo}}$ bounds the
dual objective of one specific $z$.

\subsection{The certified anchors}

Table~\ref{tab:anchors} gives, at each of the twelve centres, the
convention previously used (solver objective less $10^{-5}$) and the
certified value, at $N = 48{,}000$.

\begin{table}[h]\centering
\begin{tabular}{lrrr}
\toprule
centre & convention & certified $p_{\mathrm{lo}}$ & gain \\
\midrule
row1 & 0.380698115 & 0.381572848708 & $+8.75\times10^{-4}$ \\
row2 & 0.380677211 & 0.381665700133 & $+9.89\times10^{-4}$ \\
row3 & 0.381042228 & 0.382537653309 & $+1.50\times10^{-3}$ \\
row4 & 0.380338114 & 0.380505566439 & $+1.67\times10^{-4}$ \\
row5 & 0.380611839 & 0.380710285038 & $+9.85\times10^{-5}$ \\
row6 & 0.380436448 & 0.380552510960 & $+1.16\times10^{-4}$ \\
row7 & 0.382003546 & 0.385119737480 & $+3.12\times10^{-3}$ \\
cde1 & 0.380346079 & 0.380447891266 & $+1.02\times10^{-4}$ \\
cde2 & 0.380377646 & 0.380525911216 & $+1.48\times10^{-4}$ \\
cde3 & 0.380302283 & 0.380413948769 & $+1.12\times10^{-4}$ \\
cde4 & 0.380606798 & 0.381353402400 & $+7.47\times10^{-4}$ \\
cde5 & 0.380393675 & 0.380795566289 & $+4.02\times10^{-4}$ \\
\bottomrule
\end{tabular}
\caption{Certified anchors. Every certified value exceeds the convention
it replaces, so imposing rigour raised the bound.}\label{tab:anchors}
\end{table}

\section{Proof of Theorems~\ref{thm:main} and \ref{thm:robust}}

The cover value at a parameter point is
$\mathrm{Cover}(\theta) = \max_c \Phi_c(\theta)$, a maximum of valid
conditional lower bounds and hence itself one. Theorem~\ref{thm:main}
follows from bounding $\mathrm{Cover}$ below on every region of a
partition of the parameter space.

\subsection{Rigorous box minima}

On a box we grid and add a Lipschitz cell-error term
$\varepsilon = L_{\max}\cdot(\text{half-diagonal})$, with $L_{\max}$ the
maximum gradient magnitude of the quadratics $\Phi_c$ on that box. Since
a minimum of a maximum of concave quadratics need not be attained at a
corner, the grid-plus-Lipschitz value is the rigorous box minimum.

\begin{remark}
The certified duals are steeper than the convention they replace
($L_{\max}$ rises from $0.149$ to $0.380$ on the residual region), so a
single grid pays more than twice the $\varepsilon$ it used to, absorbing
about half the gain. Adaptive subdivision, recomputing $L_{\max}$ per
sub-box, drives $\varepsilon$ to $4.1\times10^{-8}$ and recovers it.
This trade---an honest certificate costing Lipschitz slack, repaid by
resolution---recurred at three separate points in this work.
\end{remark}

On the residual region \eqref{eq:516} this gives
$\mathrm{Cover} \ge 0.3803953504$, attained near
$(h,p) = (0.00281, 0.39219)$. The remaining regions give
$0.3803961$, $0.3803972$, $0.3803979$, $0.3804601$ and $0.3805539$, all
above it, so the residual region binds and Theorem~\ref{thm:main}
follows. Margins on the three tightest are $+7\times10^{-7}$,
$+1.8\times10^{-6}$ and $+2.5\times10^{-6}$: any future gain on the
residual region will meet them immediately.

\subsection{The robust tier}

White's Table~2 is headed ``Optimum lower bound'' and he uses its entries
as such, writing that the first line shows ``either $\mu \ge 0.38$ or
$E(M^*) \le 0.75$''. Taking his published bounds for the seventeen
regions outside the residual set, and supplying only our own certified
cover for the remaining strip and the residual region, gives
Theorem~\ref{thm:robust}. Nothing else is load-bearing.

\section{The upper bound, and a correction to the record}\label{sec:ub}

\subsection{A certified upper bound}

Any admissible step function gives an upper bound, but the published
values are float64 evaluations of float64 iterates. Every construction we
examined fails exact feasibility, by between $10^{-16}$ and
$4\times10^{-14}$ in mass---harmless in size, but enough that the stated
bounds were not theorems as written.

We snap a construction to an exactly feasible rational point and evaluate
all $2n-1$ signed lags in integer arithmetic. Applied to the best
publicly downloadable witness, after a local polish, this gives
Theorem~\ref{thm:ub}:
\[
  \mu \le \frac{129599621248162196650621937272674201309}
               {340282366920938463463374607431768211456},
\]
whose decimal expansion, rounded upward as an upper bound must be, is the
value stated. No floating-point step enters the chain.

\subsection{The value \texorpdfstring{$0.380856$}{0.380856} is not a bound}

A figure of $0.380856$ is currently quoted as the best known upper bound.
It is not one. The construction it derives from is genuinely admissible,
but its objective is $0.3809490$; the reported value is recovered
bit-for-bit only by dividing by a cell count of $4096.9999999999$ rather
than $4096$, which is what the generating program does. The honest value
is worse than constructions from four years earlier. The authors
identified and repaired this in their own repository; the paper appears
not to have been revised, and the figure continues to propagate. We note
it here because a lower-bound paper that quotes a false upper bound
misstates the size of the open gap.

\section{What this architecture cannot do}\label{sec:ceiling}

An earlier version of this paper asserted an explicit saturation ceiling
of $0.380713$. That derivation does not hold: its own per-row table
yields a value exceeding the certified upper bound, and the published
figure substitutes a min-over-rows quantity for the theorem's per-row
right-hand side. We withdraw it and replace it with the following, which
needs no sensitivity linearisation.

\begin{proposition}\label{prop:containment}
For every $N$, the value of the discretised program is at most the value
of the continuum program at the same $(T,R,\text{Bochner level})$.
\end{proposition}

Numerically, at the parameters of the extremal function, that limit is
$\approx 0.38065$, pinned two independent ways (Richardson extrapolation
on an $N$-ladder, and complementary slackness with the entire
cell-envelope slack removed). Since the certified upper bound is
$0.3808591$, this architecture is at least $2.1\times10^{-4}$ short at
the single parameter point where the extremal function lives, before any
loss from the cover. It cannot close the gap.

Two quantitative corollaries sharpen where the remaining slack is not.
An exact moment-body program---which dominates the present program plus
\emph{any} further constraint family on the $f$-side---exceeds the
certified anchor by only $3.77\times10^{-5}$ at the binding centre, and
that headroom \emph{shrinks} as $N$ grows. And the discretisation
direction obeys a clean $1/N$ law,
$\mathrm{val}(N) \approx 0.380541649 - 3.9724/N$, leaving
$8.3\times10^{-5}$ unrealised as $N \to \infty$.

\section{Approaches ruled out}\label{sec:negative}

We record five, each with a proof rather than a failed experiment.

\paragraph{Discretisation-error bounds.} One would like
$\mu \ge M_n - \varepsilon_n$ with $M_n$ the $n$-cell minimum. The
cell-average projection error satisfies
$C_{P_n h}(md) = C_h(md) + A_e(md)$ at lattice lags, and
$m \mapsto A_e(md)$ is a \emph{positive-definite} sequence, so its
maximum is at $m=0$ and equals the entire detail energy. Positivity thus
forces the error to be maximally adverse rather than allowing it to be
signed. The sharpest lattice-only bound is $\le 0$ for every $n$.

\paragraph{Single-weight lag averaging.} Replacing $\sup_k$ by
$\int w(k)\,dk$ makes the objective concave in $h$ when $\hat w \ge 0$,
so the infimum sits at an extreme point and the extreme points are
indicators. The reduction is correct and vacuous: the resulting bound is
capped at $0.3294738$ for any single probability weight and at exactly
$1/4$ for positive-definite ones. Collapsing the supremum costs at least
$0.05$, while the extreme-point structure gained is free information.

\paragraph{Argmax disjunctions.} If $P \subseteq \bigcup_j D_j$ then
$\min_j \inf\{f : P \cap D_j\} = \inf\{f : P\}$. Branching on which lag
attains the maximum is therefore identically worthless.

\paragraph{Sharpened cell envelopes.} Replacing the cell-minimum bound by
the sharp mass-and-cap (bathtub) bound is valid and strictly dominant,
but is contained in the same continuum program: it is a Richardson
extrapolator for the discretisation, obtainable more cheaply by
increasing $N$.

\paragraph{Moment-body cuts on the $f$-side.} Valid, genuinely outside
the Bochner cone, and capped at $3.8\times10^{-5}$ by the meter above.

\section{Reproducibility}\label{sec:repro}

All code and data are available \cite{repo}. Every headline number in
this paper was regenerated from scratch in a clean environment; the
full-space driver reproduces its stored result byte-identically. Two
independent re-implementations of the cover extension agree to more than
ten significant digits. We note explicitly that the helper historically
used in this project to extract ``rigorous'' bounds from solver logs is
\emph{not} a certificate---it applies no correction for dual
infeasibility, and its eligibility test read the primal residual---and
that no number in this paper depends on it.

\begin{thebibliography}{9}
\bibitem{White2023} E.~P.~White, \emph{A new bound for Erd\H{o}s'
  minimum overlap problem}, Acta Arith. \textbf{208} (2023), 235--255.
  arXiv:2201.05704.
\bibitem{KimPilanci2026} S.~Kim and M.~Pilanci, \emph{AI-assisted
  discovery of convex relaxations via dual agents}, ICML 2026, PMLR 306.
  arXiv:2606.31182.
\bibitem{JCK2007} C.~Jansson, D.~Chaykin and C.~Keil, \emph{Rigorous
  error bounds for the optimal value in semidefinite programming}, SIAM
  J. Numer. Anal. \textbf{46} (2007), 180--200.
\bibitem{Haugland2016} J.~K.~Haugland, \emph{A new upper bound on the
  constant in the Erd\H{o}s minimum overlap problem}, 2016.
  arXiv:1609.08000.
\bibitem{repo} B.~Zanghi, \emph{Erd\H{o}s minimum overlap: code and
  certificates}, \url{https://github.com/bzanghi/erdos-minimum-overlap-bochner}.
\end{thebibliography}

\end{document}
